\section{Convex specifications and optimal uniform approximation}
\label{sec:r76-convex-specifications}

The exact normal form has a robust extension. The target need not be an
exactly specified probability table, and the stochastic controller need
not satisfy the group relations, even on its observable states. What is
required is a constraint on every word representing a given group element.
The quantifier over all words is essential.

Fix a finite command alphabet $A$, a finite group $G$, and a surjective
homomorphism $\gamma:A^*\to G$. For every $g\in G$ let $K_g$ be a closed
convex subset of the finite probability simplex $\Delta(C)$. Empty sets
are allowed. An $M$-label pure-start controller consists of a prescribed
point row $e_o$, stochastic command matrices $Q_a$, and a row-stochastic
decoder $D_0$. It satisfies the specification if
\begin{equation}\label{eq:r76-specification}
 e_oQ_wD_0\in K_{\gamma(w)}\qquad(w\in A^*).
\end{equation}
For a mixed start, replace $e_o$ by an arbitrary probability row and do
not charge an additional label for preparing that row.

\begin{theorem}[Finite-group convex-specification normal form]
\label{thm:r76-convex-normal-form}
The specification \eqref{eq:r76-specification} has a pure-start
realization on at most $M$ labels if and only if there are a finite
$G$-set $X$, a probability row $\lambda$ on $X$, and a row-stochastic
decoder $D$ such that
\begin{equation}\label{eq:r76-convex-core}
 \lambda P_gD\in K_g\quad(g\in G),\qquad
 |X|+\mathbf1_{\{\lambda\text{ is not a point mass}\}}\le M.
\end{equation}
For mixed starts the corresponding criterion is $|X|\le M$.
The action on $X$ is permitted to have a kernel. In particular no
faithfulness hypothesis is imposed on an approximating action.
\end{theorem}
\begin{proof}
We give the projection argument explicitly because its use of convexity
replaces the exact observable intertwining in
Theorem~\ref{thm:r73-normal-form}.
Let $\mathcal L=\gamma^{-1}(e)$. Select finitely many identity words
$\ell_1,\ldots,\ell_s$ whose positive matrix entries cover every
directed edge occurring in any $Q_\ell$, $\ell\in\mathcal L$, and put
\[
 T=(I+Q_{\ell_1}+\cdots+Q_{\ell_s})/(s+1).
\]
At most $M^2$ words suffice. Each closed communicating class $C_i$ of
$T$ is closed under every identity-word matrix. The positive diagonal
and Lemma~\ref{lem:r73-lazy-convergence} give a stochastic idempotent
$E=\lim_nT^n$. Its image simplex has disjointly supported vertices
$\mu_1,\ldots,\mu_k$, the stationary rows of the closed classes.
Every identity word preserves those classes, and consequently
\begin{equation}\label{eq:r76-sandwich-identity}
 EQ_\ell E=E\qquad(\ell\in\mathcal L).
\end{equation}
For a word $u$ write $F_u=EQ_uE$ on this image simplex. Choose an
inverse word $\bar u$ with $\gamma(\bar u)=\gamma(u)^{-1}$.
Since $T^n$ is a convex combination of identity-word matrices,
\eqref{eq:r76-sandwich-identity} applied to $u\ell\bar u$ and
$\bar u\ell u$ gives
\[
 F_uF_{\bar u}=E=F_{\bar u}F_u.
\]
Thus $F_u$ is an affine automorphism of the image simplex and permutes
its vertices. If $\gamma(u)=\gamma(v)$, multiplication by the common
inverse $F_{\bar u}$ shows $F_u=F_v$. Replacing the middle $E$ by
$T^n$ in $EQ_uEQ_vE$ also gives $F_uF_v=F_{uv}$, since all words
$u\ell v$ have the same group image. This defines a $G$-action $P_g$
on $X=\{1,\ldots,k\}$.

Write $e_oE=\sum_i\lambda_i\mu_i$ and put $D(i,\cdot)=\mu_iD_0$.
For any representative $u$ of $g$,
\[
 \lambda P_gD=e_oEQ_uED_0
             =\lim_{n\to\infty}e_oT^nQ_uT^nD_0.
\]
The last expression before passage to the limit is a convex combination
of $e_oQ_{\ell u\ell'}D_0$ with $\ell,\ell'\in\mathcal L$.
All these probability rows belong to $K_g$. Closedness and convexity
therefore give $\lambda P_gD\in K_g$. This argument includes $u$
empty and makes no claim that the original controller's word laws
factor through $G$.

If the initial label lies in a closed class $C_i$, then $e_oE=\mu_i$
and $\lambda$ is a point mass. Otherwise the $k$ disjoint nonempty
closed classes exclude at least the initial label, so $k+1\le M$.
This proves the pure-start charge. For a mixed row $\eta$, use
$\eta E$ instead and retain $k\le M$.

Conversely, a point preparation on $X$ needs only the permutation
transitions. For a mixed $\lambda$, a further label $\star$ has decoder
$\lambda D$ and, on the first actual command $a$, moves into $X$ with
row $\lambda P_{\gamma(a)}$. Subsequent transitions are permutations.
Its empty and nonempty word laws are respectively $\lambda D$ and
$\lambda P_{\gamma(w)}D$. This gives the required realization and
finishes both start conventions.
\end{proof}

For an exact probability table $b$, let $\kappa_G(b)$ be the least
$|X|$ among factorizations $b(g)=\lambda P_gD$. This is its mixed-start
minimum by the preceding theorem. The notation $m(b)$ retains the
counted pure-start convention; its exact formula and the residual
cardinality $N$ are given in Theorem~\ref{thm:r73-memory-formula}.

For a target $b:G\to\Delta(C)$ write
$\operatorname{TV}(p,q)=\tfrac12\sum_c|p_c-q_c|$ and define
\[
 \mathcal E_M^{\mathrm p}(b)=
 \inf_{\substack{\text{pure-start controllers on at most }M\text{ labels}}}
 \ \sup_{w\in A^*}\operatorname{TV}
           (b(\gamma(w)),e_oQ_wD_0).
\]
Use $\mathcal E_M^{\mathrm m}$ for the mixed-start convention.
For a finite $G$-set $X$ put
\[
 d_X(b)=\min_{\lambda,D}\max_{g\in G}
          \operatorname{TV}(b(g),\lambda P_gD),\qquad
 d_X^{\mathrm{pt}}(b)=\min_{x,D}\max_{g\in G}
          \operatorname{TV}(b(g),e_xP_gD).
\]
The decoder ranges over stochastic rows; compactness makes both minima
attained. Minima over no candidate actions are interpreted as $+\infty$.

\begin{theorem}[Attained finite optimization for infinite-word error]
\label{thm:r76-minimax-normal-form}
For every target $b$, finite $G$, and integer $M\ge1$,
\begin{align}
 \mathcal E_M^{\mathrm m}(b)&=\min_{|X|\le M}d_X(b),\label{eq:r76-mixed-error}\\
 \mathcal E_M^{\mathrm p}(b)&=
 \min\left\{\min_{|X|\le M}d_X^{\mathrm{pt}}(b),
             \min_{|X|\le M-1}d_X(b)\right\}.
 \label{eq:r76-pure-error}
\end{align}
One representative of each finite $G$-set is sufficient, so these are
finite optimizations and their minima are attained. Their values are
$1$-Lipschitz functions of $b$ in $\max_g\operatorname{TV}$.
If $\mathcal E_M^{\mathrm p}(b)>0$, a sufficiently small perturbation of
$b$ still cannot be exactly realized on $M$ pure-start labels; the same
assertion holds for mixed starts. When the target table is algebraic,
the optimum and a realizing finite controller are computable by real
quantifier elimination. This is a decidability assertion, not a
polynomial-time bound or a claim that quantifier elimination was executed.
\end{theorem}
\begin{proof}
Apply Theorem~\ref{thm:r76-convex-normal-form} to the closed convex
balls $K_g=\{p:\operatorname{TV}(p,b(g))\le\delta\}$.
It sends every controller of error at most $\delta$ to an admissible
action without increasing the label charge or the error. Conversely
every displayed action supplies such a controller. The pure constraint
splits precisely into a point preparation with $|X|\le M$ or an
arbitrary preparation with $|X|\le M-1$.
There are finitely many homomorphisms from the fixed finite group into
each symmetric group of degree at most $M$. For each action the allowed
preparations and decoders form compact sets. Hence the finite minima
are attained, including limits in which a mixed preparation becomes a
point mass. The triangle inequality proves the Lipschitz bound and
the perturbation assertion.

For the algorithmic statement, enumerate these actions and introduce
variables $\lambda_x,D_{xc},\delta$ and $u_{gc}\ge0$ with
\[
 -u_{gc}\le b(g,c)-(\lambda P_gD)_c\le u_{gc},\qquad
 \sum_cu_{gc}\le2\delta.
\]
Together with the stochastic constraints these are polynomial
inequalities over the real algebraic numbers. The point-start branch
fixes $\lambda=e_x$ and enumerates $x$. The compact attained minimum
is the least real point of the resulting semialgebraic set of feasible
$\delta$'s. Real quantifier elimination and algebraic sample-point
construction, used with the same scope as
Theorem~\ref{thm:r73-algebraic-synthesis}, compute the value and tables.
\end{proof}

The finite optimization concerns synthesis against a supplied finite-group
target. It does not solve the cutpoint, equivalence, or isolation problem
for a supplied arbitrary probabilistic automaton. Nor does it replace
the specification for all words by a check of one representative per group
element on the original stochastic controller. Those restrictions would
change the problem.

\begin{theorem}[Robust multiplicity and initialization certificates]
\label{thm:r76-robust-multiplicity}
Use the stacked Fourier matrix $A_\rho(b)$ of
\eqref{eq:r75-fourier-block}, with singular values in decreasing order,
and write $m_\rho(X)$ for the multiplicity of $\rho$ in
$\mathbb C[X]$. If an action $X$ approximates $b$ to total-variation
error at most $\delta$, then
\begin{equation}\label{eq:r76-fourier-robust}
 \sigma_{m_\rho(X)+1}(A_\rho(b))\le\sqrt2\,\delta
                       \qquad(\rho\in\widehat G).
\end{equation}
A singular value past the number of columns means zero. For a binary
response $b=(\beta,1-\beta)$ the sharper equivalent scalar certificate is
\begin{equation}\label{eq:r76-binary-robust}
 \sigma_{m_\rho(X)+1}(\widehat\beta(\rho))\le\delta.
\end{equation}
In particular let $r_\rho^\delta$ be the number of singular values of
$A_\rho(b)$ strictly exceeding $\sqrt2\delta$. Then every mixed-start
controller of uniform error at most $\delta$ uses at least
\begin{equation}\label{eq:r76-robust-cover}
 \min\left\{\sum_H q_H[G:H]:q_H\in\mathbb Z_{\ge0},\quad
       \sum_Hq_H\ge1,\quad
       \sum_Hq_Hm_\rho(H)\ge r_\rho^\delta\ (\rho\in\widehat G)
       \right\}
\end{equation}
labels. The scalar version may be used for binary outputs.

Let $N$ be the number of distinct residual tables $b^h(g)=b(hg)$.
If $N>1$, let
\[
 \Delta_b=\min_{b^h\ne b^k}\max_g
                \operatorname{TV}(b(hg),b(kg))>0.
\]
Whenever $\delta<\Delta_b/2$, a pure-start controller with fewer than
$N$ labels must have a projected mixed preparation. Its total budget
is therefore at least one plus the lower bound
\eqref{eq:r76-robust-cover}.
\end{theorem}
\begin{proof}
Put $a(g)=\lambda P_gD$ and $z_c(g)=b(g,c)-a(g,c)$.
Since $\sum_cz_c(g)=0$ and $\tfrac12\sum_c|z_c(g)|\le\delta$,
the sums of the positive and negative parts are each at most $\delta$,
and $\sum_c|z_c(g)|^2\le2\delta^2$. Unitarity gives
\[
 \left\|\begin{pmatrix}z_{c_1}(g)\rho(g^{-1})\\
                  \vdots\\z_{c_{|C|}}(g)\rho(g^{-1})\end{pmatrix}
       \right\|_{\rm op}
 =\left(\sum_c|z_c(g)|^2\right)^{1/2}\le\sqrt2\delta.
\]
Average over $g$ to obtain
$\|A_\rho(b)-A_\rho(a)\|_{\rm op}\le\sqrt2\delta$.
The Fourier-multiplicity lemma gives
$\operatorname{rank}A_\rho(a)\le m_\rho(X)$.
For completeness, if $B$ has rank at most $m$, its nullspace
intersects the span of the first $m+1$ right singular vectors of $A$
in a nonzero vector. On a unit vector in that intersection,
$\|(A-B)v\|=\|Av\|\ge\sigma_{m+1}(A)$.
This proves \eqref{eq:r76-fourier-robust}; the binary version follows
by averaging the scalar error, whose absolute value is at most
$\delta$, against unitary matrices. Summing orbit multiplicities and
using the mixed-start normal form gives the integer bound.

For the last assertion, suppose the projected preparation were a point
$x$. Its distinct residuals are indexed by points $x\cdot h$ in its
orbit. If two target residuals $b^h,b^k$ were represented by the same
point, the two approximating residuals would coincide for every future
$g$, and the triangle inequality would give distance at most $2\delta$
between $b^h$ and $b^k$. This contradicts the definition of $\Delta_b$.
The orbit must therefore contain at least $N$ points. With fewer than
$N$ total labels this is impossible, so the projected start is mixed
and its extra initializer is charged by the normal-form theorem.
\end{proof}

\begin{corollary}[Uniform error cannot replace a prime orbit]
\label{cor:r76-prime-error}
Let $G=C_p$ with $p$ prime and let $b:G\to\Delta(C)$ be any target.
Put $r(b)=\min_{q\in\Delta(C)}\max_g\operatorname{TV}(b(g),q)$.
For either start convention,
\[
 \mathcal E_M(b)=
 \begin{cases}r(b),&1\le M<p,\\0,&M\ge p.\end{cases}
\]
For binary outputs $r(b)=(\max_g\beta(g)-\min_g\beta(g))/2$.
\end{corollary}
\begin{proof}
Every orbit of $C_p$ has size one or $p$. Thus any action with fewer
than $p$ points is trivial and gives a constant response, regardless
of preparation. A one-label controller attains $r(b)$. On $p$ points
the regular action with a point preparation and decoder $b$ is exact.
For binary outputs the midpoint of the smallest interval containing
all values of $\beta$ is a Chebyshev center of that interval.
\end{proof}

For the symmetric-group target \eqref{eq:r75-symmetric-return}, the
scalar standard block is $tI/[n(n-1)]$ and $\Delta_{b_t}=2t/n$.
Indeed a permutation moving $k$ points changes a residual by at most
$kt/n$, with equality at an inverse representative; the smallest
nonzero $k$ is two. The subgroup-index estimate
\eqref{eq:r75-index-estimate} and the preceding theorem imply
\begin{equation}\label{eq:r76-symmetric-robust-lower}
 \mathcal E_M^{\mathrm p}(b_t)\ge\frac{t}{n(n-1)}
       \qquad(n\ge5,\ 0<t\le1,\ M\le n(n-1)).
\end{equation}
If the error were smaller, the full standard multiplicity $n-1$
would be forced, and the small-error initialization charge would give
at least $n(n-1)+1$ labels; here $n!>n(n-1)$.
This explicit lower bound is not asserted to be the optimal error
constant. At every exact resource jump the finite optimization also
gives a strictly positive robustness radius, without a claimed closed
formula for that radius.
